\section{The collapsing nodal model}\label{sec:local-models}

We construct a holomorphic elliptic fibration over a disk with one
ordinary node, together with a fixed holomorphic symplectic form and
K\"ahler forms of decreasing fiber area.  This local construction
supplies the model used in the insertion argument of
\Cref{lem:pencil-insertion}; it starts independently of the global
pencil.  We use the Gibbons--Hawking ansatz
\cite{GibbonsHawking1978} in the periodic form of Ooguri--Vafa
\cite{OoguriVafa1996}, with the collapsing and semi-flat comparison
developed by Gross--Wilson \cite[Section~3]{GrossWilson2000}.

There are two further requirements for its later use.  The holomorphic
family and its symplectic form must stay fixed as the area decreases,
and the local K\"ahler form must have the same two annular periods as
its translation-invariant approximation.  We first normalize the
periodic potential and complete its node.  We then identify the
completed complex families and prove the annular comparison and
intrinsic metric estimates needed in \Cref{sec:auxiliary-volume}.

Write \(\Delta_R=\{s\in\C:|s|<R\}\), where \(0<R<1\), and set
\begin{equation}\label{eq:nodal-period}
 \tau(s)=\frac{1}{2\pi i}\log s,\qquad
 \tau_2(s)=\operatorname{Im}\tau(s)=-\frac{1}{2\pi}\log|s|.
\end{equation}
The period \(\tau\) is understood on branches, with its integral monodromy.

\begin{theorem}[The nodal model]\label{thm:nodal-model}
There are a smooth complex surface \(Y\), a proper holomorphic map
\(p:Y\to\Delta_R\), and a holomorphic section disjoint from its critical
point, with the following properties.  The central fiber is a reduced
irreducible rational curve with one ordinary node, and every other
fiber is a smooth elliptic curve.  On the regular part there are
coordinates, relative to the section,
\begin{equation}\label{eq:nodal-elliptic-coordinates}
 z=a+\tau(s)b\pmod{\Z+\tau(s)\Z},\qquad
 (a,b)\in\R^2/\Z^2,
\end{equation}
and \(ds\wedge dz\) extends as a nowhere-zero holomorphic two-form.
The real group \(H_2(Y;\R)\) is generated by a regular fiber.

Fix \(h\in\C\setminus\{0\}\), put \(\Sigma=h\,ds\wedge dz\), and fix
\(0<r<R\).  For every sufficiently small \(\epsilon>0\) there is a
K\"ahler form \(D^{\mathrm{loc}}_\epsilon\) on \(Y\), with metric
\(g_\epsilon\), such that
\begin{equation}\label{eq:nodal-exact-identities}
 \begin{gathered}
 D^{\mathrm{loc}}_\epsilon\wedge\operatorname{Re}\Sigma
 =D^{\mathrm{loc}}_\epsilon\wedge\operatorname{Im}\Sigma=0,\qquad
 (D^{\mathrm{loc}}_\epsilon)^2=(\operatorname{Re}\Sigma)^2,\\
 \int_{p^{-1}(s)}D^{\mathrm{loc}}_\epsilon=\epsilon\quad(s\ne0).
 \end{gathered}
\end{equation}
These three forms are parallel, mutually orthogonal and of equal
square, and \(\vol_{g_\epsilon}=(\operatorname{Re}\Sigma)^2/2\).
Moreover:
\begin{enumerate}[label=\textup{(\roman*)}]
\item On each fixed closed annulus
\(K=\{r_-\le |s|\le r_+\}\subset\Delta_R\setminus\{0\}\), the form
\begin{equation}\label{eq:nodal-semiflat}
 D^{\mathrm{sf}}_\epsilon
   =\epsilon\,da\wedge db+
      \frac{|h|^2}{\epsilon}\tau_2(s)\,dx_1\wedge dx_2,
       \qquad s=x_1+ix_2,
\end{equation}
is well-defined and closed.  For every integer \(m\ge0\), in fixed
smooth charts on this annular bundle,
\begin{equation}\label{eq:nodal-exponential-estimate}
 \|D^{\mathrm{loc}}_\epsilon-D^{\mathrm{sf}}_\epsilon\|_{C^m}
       \le C_m e^{-c_m/\epsilon},
       \qquad C_m,c_m>0.
\end{equation}
\item These forms have the same real cohomology class on \(p^{-1}(K)\).
Their periods are \(\epsilon\) on a fiber and zero on the torus
sweeping the invariant \(a\)-circle at \(b=0\) around the base annulus.
That swept torus bounds in \(Y\).
\item On \(p^{-1}(\overline{\Delta_r})\), for every fixed \(m\) there
is a finite coordinate cover with smaller concentric domains still
covering this set, such that the number of charts, the reciprocals
of their radii and interior margins, and the \(C^m\) bounds for
the metric coefficients and their inverses are at most
\(C_m\epsilon^{-N_m}\).  All intrinsic curvature derivatives of
fixed finite order have polynomial bounds as well.
\end{enumerate}
All disks, annuli and \(h\) are fixed before \(\epsilon\) tends to zero.
The charts in \textup{(iii)} may depend on \(\epsilon\).
\end{theorem}

Thus, for fixed real \(x_0,y_0\), \(y_0\ne0\), the pair
\[
 A=\operatorname{Re}\Sigma,\qquad
 C=x_0\operatorname{Re}\Sigma+y_0\operatorname{Im}\Sigma
\]
is orthogonal to \(D^{\mathrm{loc}}_\epsilon\).  Its span with this
form is positive, and its metric in volume \(A^2/2\) is \(g_\epsilon\).
No normalization of \(A,C\) relative to one another is required.

\subsection{The periodic potential}

The normalization follows the periodic-potential construction in
\cite[Lemma~3.1 and Proposition~3.2]{GrossWilson2000}.
First take \(h=1\), and denote the area parameter by \(\rho\).
On \((\Delta_R\times\R/\rho\Z)\setminus\{(0,0)\}\) use coordinates
\((x_1,x_2,t)\).
For a precise choice of the additive constant define
\begin{equation}\label{eq:nodal-unit-potential}
\begin{split}
 \mathcal V(u,v)=\frac{1}{4\pi}\bigg\{
 &\frac{1}{\sqrt{|v|^2+u^2}}\\
 &+\sum_{n=1}^\infty
 \left(\frac{1}{\sqrt{|v|^2+(u+n)^2}}+
       \frac{1}{\sqrt{|v|^2+(u-n)^2}}-\frac2n\right)\bigg\}.
\end{split}
\end{equation}
The paired summands are \(O(n^{-3})\) locally, with locally uniformly
convergent differentiated series away from the poles.  The resulting
function is smooth, harmonic, and periodic with period one in \(u\).
Its mean is
\begin{equation}\label{eq:nodal-unit-mean}
 \int_{-1/2}^{1/2}\mathcal V(u,v)\,du
  =-\frac{\log|v|}{2\pi}
       +\frac{\log2-\gamma_{\mathrm E}}{2\pi},
\end{equation}
where \(\gamma_{\mathrm E}\) is Euler's constant.  Indeed, the integral
of the sum truncated at \(|n|\le N\), before the factor \(1/(4\pi)\),
is
\[
 2\operatorname{arsinh}\frac{N+1/2}{|v|}
      -2\sum_{n=1}^N\frac1n
 \longrightarrow 2\log(2/|v|)-2\gamma_{\mathrm E}.
\]
Consequently
\begin{equation}\label{eq:nodal-scaled-potential}
 V_\rho(s,t)=\frac1\rho
 \left\{\mathcal V\left(\frac t\rho,\frac s\rho\right)
       -\frac{\log\rho}{2\pi}
       +\frac{\gamma_{\mathrm E}-\log2}{2\pi}\right\}
\end{equation}
has a positive unit monopole \(1/(4\pi\sqrt{|s|^2+t^2})\), and
\begin{equation}\label{eq:nodal-mean}
 \frac1\rho\int_0^\rho V_\rho(s,t)\,dt
       =\frac{\tau_2(s)}{\rho}.
\end{equation}
The term \(-(\log\rho)/(2\pi)\) in the scaled potential is essential
for keeping the periods fixed.

For \(s\ne0\), its Fourier expansion is
\begin{equation}\label{eq:nodal-fourier}
 V_\rho(s,t)=\frac{\tau_2(s)}{\rho}
  +\frac1{2\pi\rho}\sum_{n\ne0}e^{2\pi int/\rho}
       K_0\left(\frac{2\pi|ns|}{\rho}\right).
\end{equation}
It follows by the Gaussian integral representation of the Newton
kernel and integration in the periodic variable.  The nonzero
Fourier coefficient is proportional to
\[
 \int_0^\infty \ell^{-1}
     \exp\left(-\ell|s/\rho|^2-\frac{\pi^2n^2}{\ell}\right)d\ell.
\]
Alternatively \(K_0(a)=\int_0^\infty e^{-a\cosh u}\,du\).
These representations, including their differentiated versions, give
for each fixed annulus and each \(m\)
\begin{equation}\label{eq:nodal-potential-tail}
 \left\|V_\rho-\frac{\tau_2}{\rho}\right\|_{C^m
  (\{r_-\le|s|\le r_+\}\times\R/\rho\Z)}
   \le C_m\rho^{-N_m}e^{-2\pi r_-/\rho}.
\end{equation}
Derivatives in \(t\) are taken on periodic local lifts.
Differentiation introduces powers of \(\rho^{-1}\) and \(|n|\);
the exponential decay makes the corresponding sums convergent.
Any smaller positive exponential rate absorbs the polynomial
prefactors.

The function \(V_\rho\) is positive on the whole punctured solid
torus when \(\rho\) is sufficiently small.  Choose \(L>1\) large
enough that the Fourier tail of \(\rho V_\rho\) on
\(|s|/\rho\ge L\) is less than \(-(\log R)/(4\pi)\).
Its mean is at least \(-(\log R)/(2\pi)>0\), proving positivity
there.  On the remaining compact scaled region \(|s|/\rho\le L\),
the function \(\mathcal V\) is bounded below, including near its
positive pole.  The additive term \(-(\log\rho)/(2\pi)\) then
proves positivity there also.  This argument gives a positive lower
bound for \(\rho V_\rho\) off a fixed scaled pole neighborhood.

\subsection{The metric and the smooth nodal fiber}

The Euclidean Hodge star uses the orientation
\(dx_1\wedge dx_2\wedge dt\).
The form \(*dV_\rho\) has integral \(-1\) on a small positively
oriented linking sphere.  That sphere generates the second homology
of the punctured solid torus.  Hence it is the curvature of a real
connection \(\theta\) on a principal \(\R/\Z\)-bundle, normalized by
integral one on its circle fibers:
\begin{equation}\label{eq:nodal-connection}
 d\theta=*dV_\rho.
\end{equation}
In this period-one convention the curvature itself represents the
first Chern class.  We may change \(\theta\) by a closed basic
one-form.

With \(x_3=t\), set, for cyclic \((i,j,k)\),
\begin{equation}\label{eq:nodal-gh-triple}
 \begin{split}
 g_\rho&=V_\rho(dx_1^2+dx_2^2+dt^2)+V_\rho^{-1}\theta^2,\\
 \gamma_i&=dx_i\wedge\theta+V_\rho\,dx_j\wedge dx_k.
 \end{split}
\end{equation}
The signs follow directly from
\[
 d(dx_i\wedge\theta)=-\partial_iV_\rho\,dx_1\wedge dx_2\wedge dx_3,
 \qquad
 d(V_\rho dx_j\wedge dx_k)=\partial_iV_\rho\,dx_1\wedge dx_2\wedge dx_3.
\]
Thus all three forms are closed.  In the coframe
\(V_\rho^{1/2}dx_i,V_\rho^{-1/2}\theta\) they form the standard
orthogonal equal-square self-dual frame, with the induced
four-dimensional orientation.  To check parallelness, write the
connection on this frame in terms of three real one-forms \(a_i\).
After a consistent choice of frame signs, the equations for
closedness are
\[
 I_2a_3-I_3a_2=0,\qquad
 I_3a_1-I_1a_3=0,\qquad
 I_1a_2-I_2a_1=0,
\]
where \(I_i\alpha=*(\alpha\wedge\gamma_i)\), with an overall
sign chosen so \(I_1I_2=I_3\).
The first two equations give \(a_3=-I_1a_2\) and
\(a_1=-I_2a_3=-I_3a_2\).
The third then gives \(2I_1a_2=0\), so all \(a_i\) vanish.
The frame is parallel.
In particular
\begin{equation}\label{eq:nodal-gh-complex-form}
 \Omega_\rho=\gamma_1+i\gamma_2
       =ds\wedge(\theta-iV_\rho dt).
\end{equation}
It is decomposable and has positive product with its conjugate.
Its kernel defines the complex structure for which \(\gamma_3\)
is K\"ahler.  This kernel is involutive: closedness and Cartan's
formula give \(\iota_{[U,W]}\Omega_\rho=0\) whenever \(U,W\)
annihilate \(\Omega_\rho\).  The function \(s\) is holomorphic.

Here is the completion and the quantitative information at its pole.
In scaled coordinates \(\mathbf X=(x_1,x_2,t)/\rho\), write
\begin{equation}\label{eq:nodal-pole-split}
 W=\rho V_\rho=W_m+H_\rho,\qquad W_m=\frac1{4\pi r},\quad
 r=|\mathbf X|.
\end{equation}
On a fixed ball of radius less than \(1/4\),
\(H_\rho=H_0-(\log\rho)/(2\pi)\), where \(H_0\) is smooth harmonic
and independent of \(\rho\).
A local gauge gives \(\theta=\theta_m+\beta\), where
\(d\theta_m=*dW_m\) and \(\beta\) is smooth basic with
\(d\beta=*dH_0\).  Its derivatives on a smaller ball are bounded
independently of \(\rho\).

The charge-one Hopf chart completes
\(W_m\,d\mathbf X^2+W_m^{-1}\theta_m^2\) to a flat smooth
four-dimensional ball.  In this chart \(\mathbf X\) and \(r\) are
smooth quadratic expressions, and \(r\theta_m\) is smooth.
The coefficient \(1/(4\pi)\), with connection period one, is the
smooth Hopf normalization.  Now
\begin{equation}\label{eq:nodal-pole-inverse}
 W^{-1}=\frac{4\pi r}{1+4\pi rH_\rho},\qquad
 \frac{W^{-1}-W_m^{-1}}{r^2}
    =-\frac{(4\pi)^2H_\rho}{1+4\pi rH_\rho}.
\end{equation}
It follows that
\[
 g_\rho/\rho=W\,d\mathbf X^2+
                   W^{-1}(\theta_m+\beta)^2
\]
is smooth at the added point.  In particular, the coefficient of
the apparently singular squared-connection difference is smooth
by the second identity in \Cref{eq:nodal-pole-inverse}.
Subtracting the flat Hopf triple in
\[
 \gamma_i/\rho=dX_i\wedge(\theta_m+\beta)+W\,dX_j\wedge dX_k
\]
shows that the forms also extend smoothly.  At the added point
their values are the flat triple, so they remain nondegenerate.
We obtain a smooth hyperk\"ahler surface \(Y_\rho\), and its
holomorphic function \(s\) extends over the added point.

At that point \(ds=0\).  The quadratic term of \(s\) is the
nondegenerate holomorphic quadratic of the Hopf map; in linear
complex coordinates it is a nonzero multiple of \(w_1w_2\).
The holomorphic Morse lemma therefore gives an ordinary node.
Away from this point the central fiber is a circle bundle over
\((\R/\rho\Z)\setminus\{0\}\).  Completing its two ends at the pole
gives a reduced irreducible nodal curve with spherical normalization.
The fibers for \(s\ne0\) are two-tori, and their relative holomorphic
one-form from \(\Omega_\rho\) has no zeros; they are elliptic curves.
The projection \(p_\rho:Y_\rho\to\Delta_R\) is proper: over a compact
base set its complement of a pole neighborhood is a circle bundle
over a compact set, and the missing part is a compact subset of
the completed Hopf chart.

The potential is invariant under \(t\mapsto\rho-t\), so the
curvature restricts to zero on \(t=\rho/2\).
The flat circle connection over this disk has a horizontal section
\(S_\rho\).  It is holomorphic: its tangent space has
\(dt=\theta=0\), so \(\Omega_\rho\) vanishes there, while \(ds\)
is an isomorphism to the base tangent.  It avoids the critical point.

Integrate the relative holomorphic differential from this section
to obtain the Abel coordinate \(z\) on each smooth fiber.
The connection circle has period one.  Orient the other cycle in
the negative \(t\)-direction.  Its period \(\tau_\rho\) then satisfies
\[
 \operatorname{Im}\tau_\rho(s)
       =\int_0^\rho V_\rho(s,t)\,dt=\tau_2(s).
\]
The periods are holomorphic on each branch of the regular base.
Thus \(\tau_\rho-\tau\) is a real constant.
Adding a real multiple of \(dt/\rho\) to \(\theta\) changes exactly
this real period; we choose it to make \(\tau_\rho=\tau\).
An integral change only changes the cycle basis.
This does not affect curvature or the horizontal section in the
midpoint slice.  In the resulting coordinates
\(\Omega_\rho=ds\wedge dz\).

\subsection{The fixed complex surface across the node}

The common periods identify the regular elliptic families,
preserving their sections and relative differentials.
The following argument extends the identification across zero.

\begin{lemma}\label{lem:nodal-fixed-family}
The preceding identifications extend to holomorphic symplectic
isomorphisms of the completed families \(Y_\rho\to\Delta_R\).
\end{lemma}
\begin{proof}
We compare the families through the degree-three divisors given by
three times their sections.  Their line bundles embed the fibers as
plane cubics.  Laurent coefficients at the section, measured in the
normalized Abel coordinate, will identify the three-dimensional spaces
of sections across the puncture; taking closures will then identify
the completed families.

Each \(p_\rho\) is a flat proper family of connected reduced curves.
For flatness, the local ring of the smooth total surface has no
torsion over the local ring of the one-dimensional smooth base,
since \(p_\rho-s_0\) is not a zero divisor.  Torsion-free modules
over this discrete valuation ring are flat.
The fibers have arithmetic genus one, and \(S_\rho\) is a Cartier
divisor in the fiberwise smooth locus.

Set \(\mathcal L_\rho=\mathcal O_{Y_\rho}(3S_\rho)\).
This line bundle is flat over the base, since it is locally free on
the flat family.
Its restriction to each fiber is a degree-three very ample line
bundle, with \(h^0=3\) and \(h^1=0\).
For the complex-linear Dolbeault operator on a line bundle of
degree \(d\) over a closed connected Riemann surface of genus \(g\), the complex
index is \(d+1-g\), and \(d>2g-2\) gives surjectivity
\cite[Section~3.4, Theorems~3.22--3.23]{Wendl2014Lectures}.
On a smooth genus-one fiber this gives the asserted dimensions.
After imposing any length-two subscheme, the remaining degree is
one and the first cohomology still vanishes.  This proves very
ampleness on those fibers.
For the nodal curve one can check the assertion explicitly.
Normalize it by \(\mathbb P^1\), with the two preimages of the
node at \(0,\infty\).  A degree-three line bundle pulls back to
\(\mathcal O_{\mathbb P^1}(3)\), with a nonzero gluing parameter
\(\lambda\) for its fibers at these points.
Its sections are degree-at-most-three polynomials with
constant coefficient \(\lambda\) times the leading coefficient.
The basis \(w,w^2,w^3+\lambda\) gives the map
\[
 w\longmapsto[w:w^2:w^3+\lambda].
\]
It identifies exactly \(0\) and \(\infty\); away from these
points the ratio of the second coordinate to the first is \(w\).
At \(0\) and \(\infty\) the two tangent lines are distinct, as
is seen in the affine chart about \([0:0:1]\).
Thus it embeds the nodal curve as a nodal cubic.  To compute its
cohomology directly, let \(N\) denote the nodal fiber,
\(\nu:\mathbb P^1\to N\) its normalization, and
\(L=\mathcal L_\rho|_N\).  The gluing condition gives the exact sequence
\[
 0\longrightarrow L\longrightarrow
 \nu_*\mathcal O_{\mathbb P^1}(3)
 \xrightarrow{\operatorname{ev}_0-\lambda\operatorname{ev}_\infty}
 \C_{\mathrm{node}}\longrightarrow0,
\]
where the last sheaf is supported at the node and its fiber is \(\C\).
In the polynomial trivializations just used, the map on global
sections is the constant coefficient minus \(\lambda\) times the
leading coefficient.  It is surjective.  The normalization is finite,
so the cohomology of its direct image is that of
\(\mathcal O_{\mathbb P^1}(3)\).  The latter has \(h^0=4\) and
\(h^1=0\).  The cohomology exact sequence therefore gives
\(h^0(N,L)=3\) and \(h^1(N,L)=0\).

We use the following proper holomorphic base-change theorem: for a
proper holomorphic map to a reduced base and a coherent sheaf flat
over that base, constancy of the fiber dimensions \(h^q\) implies
that the \(q\)-th direct image is locally free and its restriction
to each base point is the cohomology of that fiber; see
\cite[Section~7, Satz~5]{Grauert1960} and the correction
\cite{Grauert1963}.  Our map is proper and
surjective, the disk is reduced, and \(\mathcal L_\rho\) is flat
over it.  Applying the theorem with \(q=0\) and the constant value
\(h^0=3\) gives a holomorphic rank-three bundle
\[
 E_\rho=(p_\rho)_*\mathcal L_\rho
\]
with those spaces as fibers.  Evaluation gives a relative closed
embedding in \(\mathbb P(E_\rho^*)\).
Indeed the fiberwise embeddings and base change give the
evaluation and separation properties locally over the base.
At the nodal point, the fiber's Zariski tangent space is the whole
tangent space of the smooth total surface, so tangent separation
there also gives an immersion of the total surface.  Properness
then gives the closed embedding.

Near the section \(p_\rho\) is a submersion.  Write
\(\Omega_\rho=f(s,w)\,ds\wedge dw\), with \(w=0\) on the section.
The coordinate
\[
 z(s,w)=\int_0^w f(s,u)\,du
\]
is holomorphic, has nonzero relative derivative, and is zero on
the section.  It extends the normalized local Abel coordinate
also across \(s=0\).
Regard sections of \(\mathcal L_\rho\) as meromorphic functions
with poles of order at most three on \(S_\rho\).
Their Laurent coefficients in degrees \(-3,-2,-1,0\) define
\begin{equation}\label{eq:nodal-jet-map}
 E_\rho\longrightarrow\mathcal O_{\Delta_R}^{\,4}.
\end{equation}
This map is fiberwise injective.  A function in its kernel has
no pole and hence is constant on the compact connected reduced
fiber; its zero constant coefficient makes it zero.
Thus its image is a rank-three subbundle, including over zero.

The images agree on the punctured disk by the regular elliptic
identifications.  Continuity of their Grassmannian-valued maps
makes them agree also at zero, identifying the bundles \(E_\rho\).
Their relative cubic images agree as well: each is the closure
of its image over the punctured disk, and both are reduced
analytic subspaces.  This gives an isomorphism of the completed
families.  The holomorphic two-forms agree on the regular part
and therefore everywhere.
\end{proof}

Fix one small reference parameter and take its surface as \(Y\).
Using \Cref{lem:nodal-fixed-family}, the form
\(\Sigma_0=ds\wedge dz\) is fixed on \(Y\), while the transported
forms \(\gamma_3\) are its varying K\"ahler forms.
In particular the qualitative model existence used to insert
nodes does not assume a global fibration.

\subsection{Comparison and periods on fixed annuli}

The complex family is now fixed.  To glue its K\"ahler forms later,
we need estimates in the same smooth annular coordinates for every
area parameter, together with equality of their annular classes.
Work over a fixed annulus with positive inner radius and on
bounded branches relative to the midpoint section.  The imaginary
part of the Abel coordinate is
\begin{equation}\label{eq:nodal-q-coordinate}
 q=\operatorname{Im}z=\int_t^{\rho/2}V_\rho(s,u)\,du.
\end{equation}
Consequently \Cref{eq:nodal-potential-tail} gives
\begin{equation}\label{eq:nodal-q-tail}
 q=\frac{\rho/2-t}{\rho}\tau_2(s)+E_\rho(s,t),\qquad
 \|E_\rho\|_{C^m}\le C_m\rho^{-N_m}e^{-c/\rho}.
\end{equation}
The interval of integration has length \(O(\rho)\); differentiated
endpoints cost only powers of \(\rho^{-1}\).
Since \(q=\tau_2b\) and \(\tau_2\) is bounded above and below,
\begin{equation}\label{eq:nodal-inverse-t}
 b=\frac{\rho/2-t}{\rho}+O(e^{-c'/\rho}),\qquad
 t=\frac{\rho}{2}-\rho b+O(e^{-c'/\rho})
\end{equation}
with all fixed finite-order derivatives in the corresponding
coordinates.  For the inverse statement,
\(\partial_tb=-V_\rho/\tau_2\) has absolute value bounded below
by a constant times \(\rho^{-1}\).  Successive differentiation
of the inverse relation introduces only polynomial factors,
absorbed by reducing \(c'>0\).

There is a complex function \(k\) on each branch such that
\begin{equation}\label{eq:nodal-connection-comparison}
 dz-(\theta-iV_\rho dt)=k\,ds.
\end{equation}
The difference gives zero relative differential and has zero
wedge with \(ds\), by \Cref{eq:nodal-gh-complex-form}.
Its imaginary part says
\[
 \operatorname{Im}(k\,ds)=d_s q
       =b\,d\tau_2+O(e^{-c'/\rho}),
\]
where \(t\) is held fixed.  Since \(\tau\) is holomorphic,
\(\operatorname{Im}(b\tau'(s)\,ds)=b\,d\tau_2\).
The real-linear map taking \(k\) to \(\operatorname{Im}(k\,ds)\)
is an isomorphism onto real base one-forms.  Thus
\(k=b\tau'(s)+O(e^{-c'/\rho})\).
Substituting \(dz=da+\tau\,db+b\tau'(s)\,ds\) in
\Cref{eq:nodal-connection-comparison} gives
\begin{equation}\label{eq:nodal-theta-tail}
 \theta=da+\operatorname{Re}\tau\,db+O(e^{-c'/\rho}).
\end{equation}
All the errors here have their fixed finite-order derivative
bounds in the fixed \((s,a,b)\) charts.  Therefore
\begin{equation}\label{eq:nodal-area-one-comparison}
 \begin{split}
 \gamma_3&=dt\wedge\theta+V_\rho dx_1\wedge dx_2\\
 &=\rho\,da\wedge db+
       \frac{\tau_2(s)}{\rho}\,dx_1\wedge dx_2
            +O(e^{-c'/\rho}).
 \end{split}
\end{equation}
The fiber sign follows from
\((-\,\rho\,db)\wedge da=\rho\,da\wedge db\).
Its exact area is \(\rho\), since on a fiber
\(\gamma_3=dt\wedge\theta\), and the two periods are \(\rho\)
and one, with this complex orientation.

Under monodromy \(\tau\mapsto\tau+1\), the angle change is
\((a,b)\mapsto(a-b,b)\).  It preserves \(da\wedge db\), and
\(\tau_2\) is single-valued.  Hence the invariant comparison
form and the finite-chart estimates descend to the annular bundle.

To identify the periods that test exactness on the annular bundle,
we first determine the topology of the completed disk neighborhood.
A one-critical-point Lefschetz fibration over a disk has the
homotopy type of a regular fiber with its vanishing thimble attached.
Here this description follows by trivializing over a slit disk:
the circle over a path to the critical value collapses to the Hopf
point and gives the one attached disk, while the complement of
that path is a trivial smooth fibration.  The vanishing circle is
the primitive connection circle \(a\).
Thus the cellular description is a torus with one additional
two-cell attached along \(a\).  Collapsing that cell together
with \(a\) leaves a one-cell \(b\) and the torus two-cell with
trivial attaching map, so the homotopy type is \(S^1\vee S^2\).
More precisely its cellular boundary sends the additional
two-cell to \(a\) and sends the torus two-cell to zero.
Consequently its second real homology is generated by the
original torus fiber, as asserted.

\begin{lemma}\label{lem:nodal-annular-periods}
The two forms in \Cref{eq:nodal-area-one-comparison} have the same
real cohomology class on the annular bundle: their periods are
\(\rho\) on a fiber and zero on the invariant-cycle sweep.
\end{lemma}
\begin{proof}
We use the two-cycle comparison from
\cite[proof of Lemma~4.3]{GrossWilson2000}.
The annular bundle retracts to the mapping torus of one shear.
Its homology exact sequence is
\[
 0\longrightarrow H_2(T^2;\R)
 \longrightarrow H_2(p^{-1}(K);\R)
 \longrightarrow\ker(T_*-1:H_1(T^2;\R)\to H_1(T^2;\R))
 \longrightarrow0.
\]
The first term is generated by a fiber, and the last by the
invariant \(a\)-circle.  This circle has a sweep at \(b=0\).
These two cycles thus test all real two-periods.

The fiber periods were computed above.  The sweep is the
connection circle bundle over a base circle in \(t=\rho/2\).
The connection circle bundle over the full base disk in this
slice is a solid torus bounding it in \(Y\).
On the sweep \(dt=0\), and the base has only one tangent
direction, so the restriction of
\(dt\wedge\theta+V_\rho dx_1\wedge dx_2\) is zero.
The invariant comparison form also has zero restriction, because
\(db=0\) and the same dimensional observation applies.
\end{proof}

\subsection{Polynomial intrinsic geometry}

The annular estimates use a fixed smooth atlas, as the gluing requires.
For the remaining metric estimates the charts may instead vary with
\(\rho\): we transport each metric together with its own controlled
charts.  This is enough to bound intrinsic derivatives on the model
cores in \Cref{sec:auxiliary-volume}.  We prove these bounds on a fixed
smaller disk.

Near the pole use the Hopf chart of \Cref{eq:nodal-pole-split}.
For small \(\rho\), the function \(H_\rho\) is positive on this
chart; its positive-order derivatives are bounded and
\(|H_\rho|\le C(1+|\log\rho|)\).
The formulas \Cref{eq:nodal-pole-inverse} show that the
coefficients of \(g_\rho/\rho\) and their derivatives have
polynomial bounds in \(1+|\log\rho|\).
For the inverse bounds compare
\[
 W\,d\mathbf X^2+W^{-1}(\theta_m+\beta)^2
 \quad\hbox{with}\quad
 W_m\,d\mathbf X^2+W_m^{-1}\theta_m^2.
\]
The ratio \(W/W_m=1+4\pi rH_\rho\) and its reciprocal have
polynomial bounds.  Replacing \(\theta_m\) by
\(\theta_m+\beta\) is a shear with smooth bounded coefficients
in the flat Hopf metric, because \(\beta\) is basic and smooth
in the quadratic base coordinates.  This gives upper and lower
polynomial metric comparisons.  Differentiation of the matrix
inverse gives its derivative estimates.  Restoring the factor
\(\rho\) yields polynomial bounds in \(\rho^{-1}\).

Off a smaller scaled pole neighborhood, use fixed-radius balls
in scaled coordinates \((s/\rho,t/\rho)\).
On such balls over \(\overline{\Delta_r}\), the periodic sum and
the Fourier formula give, for every multi-index,
\begin{equation}\label{eq:nodal-scaled-regular-bounds}
 c\le W,\qquad W\le C(1+|\log\rho|),\qquad
 |\partial^\alpha W|\le C_\alpha(1+|\log\rho|).
\end{equation}
The lower bound follows from positivity on the fixed larger
disk.  Bounded scaled distance uses the periodic sum; large
scaled distance uses the mean and Fourier tails.
The curvature is \(*dW\) in these coordinates, so radial local
gauges give connection coefficients with the same finite-order
bounds.  Cover the period-one connection circle by fixed
coordinate intervals.  The metric in these charts is
\[
 \rho\{W\,d\mathbf X^2+
              W^{-1}(d\vartheta+\beta)^2\}.
\]
\Cref{eq:nodal-scaled-regular-bounds} gives polynomial bounds for
its coefficients, inverse and all fixed derivatives.

There are polynomially many charts with the required margins.
The scaled base disk has radius \(O(\rho^{-1})\), and its
periodic direction has length one.  A grid of fixed-radius balls,
with a fixed number of circle charts above each, uses at most
\(C\rho^{-2}\) charts off the pole.  Refining the grid by a
fixed factor makes smaller concentric balls still cover.
Use similarly smaller circle intervals.  The fixed gap \(R-r\)
supplies room at the outer boundary, and the Hopf chart covers
the omitted pole region.  Coordinate formulas for the connection
and curvature, with the inverse metric estimates, also give
polynomial bounds for each fixed intrinsic curvature derivative.

These bounds survive transport to the fixed complex surface.
If \(I_\rho:Y\to Y_\rho\) is the isomorphism from
\Cref{lem:nodal-fixed-family}, transport the charts by
\(I_\rho^{-1}\) and the metric by \(I_\rho^*\).
Their metric coefficients are exactly the coefficients already
estimated, and
\[
 I_\rho^*(\gamma_1+i\gamma_2)=\Sigma_0.
\]
Thus the fixed real and imaginary parts of \(\Sigma_0\) and the
transported \(\gamma_3\) are parallel, and the volume is the
fixed form \((\operatorname{Re}\Sigma_0)^2/2\).
No bound for derivatives of \(I_\rho\) in a separately fixed
smooth atlas is used.  On annuli, where a fixed atlas is needed
for gluing, \Cref{eq:nodal-q-coordinate,eq:nodal-theta-tail}
give that distinct estimate.

\begin{proof}[Completion of the proof of \Cref{thm:nodal-model}]
The construction and \Cref{lem:nodal-fixed-family} give the
fixed holomorphic nodal disk.  For general \(h=|h|e^{i\phi}\),
rotate \(\gamma_1+i\gamma_2\) by \(e^{i\phi}\), scale the full
triple and the metric by \(|h|\), and take
\[
 \rho=\epsilon/|h|,\qquad
 D^{\mathrm{loc}}_\epsilon=|h|\gamma_3,\qquad
 g_\epsilon=|h|g_\rho.
\]
The holomorphic two-form becomes \(\Sigma=h\,ds\wedge dz\), the fiber area
becomes \(\epsilon\), and the equal-square identities give
\Cref{eq:nodal-exact-identities}.  The scaled annular expression
is precisely \Cref{eq:nodal-semiflat}, with coefficient
\(|h|^2/\epsilon\).  The exponential comparison and both period
identities follow from
\Cref{eq:nodal-area-one-comparison,lem:nodal-annular-periods}.
The polynomial estimates persist because \(h\) is fixed and
nonzero.  This proves all assertions.
\end{proof}
